\documentclass[10pt,twocolumn]{article}

% ── Packages ──────────────────────────────────────────────────────────────────
\usepackage[margin=1in]{geometry}
\usepackage{amsmath,amssymb,amsthm}
\usepackage{graphicx}
\usepackage{booktabs}
\usepackage{multirow}
\usepackage{array}
\usepackage{xcolor}
\usepackage{hyperref}
\usepackage{algorithm}
\usepackage{algpseudocode}
\usepackage{microtype}
\usepackage{natbib}
\usepackage{caption}
\usepackage{subcaption}

\hypersetup{colorlinks=true, linkcolor=blue, citecolor=blue, urlcolor=blue}

% ── Title ─────────────────────────────────────────────────────────────────────
\title{\textbf{DeepPySR: Interpretable Scientific Discovery via Dynamic
Feature Pruning and Multi-Layer Symbolic Regression}}

\author{
  Fuling Chen$^{1}$\thanks{Correspondence: \texttt{fuling.chen@icrar.org}} \\
  $^1$ International Centre for Radio Astronomy Research (ICRAR),
       University of Western Australia
}

\date{}

% ── Document ──────────────────────────────────────────────────────────────────
\begin{document}
\maketitle

% ─────────────────────────────────────────────────────────────────────────────
\begin{abstract}
Symbolic regression (SR) discovers closed-form mathematical equations from
data—enabling human-interpretable models in science and medicine.  Despite
rapid progress, two persistent challenges limit its applicability: (i)
\emph{high-dimensional feature spaces} cause search inefficiency as the
algorithm explores irrelevant variables, and (ii) \emph{formula selection}
from the Pareto front of accuracy–complexity trade-offs remains largely
heuristic.  We introduce \textbf{DeepPySR}, an enhanced SR framework built
on SymbolicRegression.jl, with three novel contributions.
\textbf{First}, we propose a \emph{dynamic variable pruning schedule} that
linearly ramps up mutation pressure to eliminate uninformative variables during
the evolutionary search, substantially reducing search space and preventing
overfitting through spurious features.
\textbf{Second}, we introduce an \emph{exponential Pareto selection criterion}
that jointly optimises R$^2$ and equation complexity via a principled
score function, replacing ad-hoc thresholding.
\textbf{Third}, we extend the framework to \emph{multi-layer symbolic
regression} (DeepSR), where intermediate latent functions are automatically
discovered and composed hierarchically, broadening the class of representable
relationships.
We evaluate DeepPySR on five Feynman physics benchmarks and five real-world
biomedical datasets.  DeepPySR recovers exact physics equations (R$^2 = 1.0$)
from noisy data and outperforms both PySR and all black-box baselines (XGBoost,
Random Forest, MLP, KAN) on body fat (R$^2$: 0.794 vs.\ 0.702), heart disease
(F1: 0.898 vs.\ 0.787), and diabetes prediction, while returning
interpretable analytical formulas that connect to domain knowledge.
\end{abstract}

% ─────────────────────────────────────────────────────────────────────────────
\section{Introduction}
\label{sec:intro}

The tension between predictive accuracy and interpretability is one of the
central problems in applied machine learning.  In scientific disciplines—from
biomedicine to astrophysics—understanding \emph{why} a model makes a
prediction is often as important as the prediction itself.  A cardiologist
needs not just a risk score but the mechanistic factors driving it; an
epidemiologist studying childhood BMI trajectories needs an equation that can
be communicated and acted on.

[comment] add the data problem. In most biomedical datasets, severe ICC or data imbalance make any feature engineering
or the available machine learning models frustrated. The available solution is to select a subset of critical
features if large amounts of variables are provided or use SMOTE alike methods to make up synthetic data to
solve the data imbalance. This is not an ideal way.

Symbolic regression (SR) addresses this need by searching for analytical
mathematical expressions that fit observed data~\citep{koza1994genetic}.
Unlike black-box models, SR outputs human-readable equations that can be
inspected, simplified, and connected to physical laws.  The field has seen
renewed interest following the release of PySR~\citep{cranmer2023interpretable}
and the Feynman SR benchmark~\citep{udrescu2020ai}, with recent methods achieving
exact recovery of fundamental physics laws from noisy measurements.

[comment] add KAN model which is a potential replacement of MLP with symbolic regression function supported.

\subsection{Limitations of existing approaches}

Despite this progress, three challenges hinder SR in practical, high-dimensional
settings:

\textbf{(1) Curse of dimensionality.}  Genetic programming (GP) explores an
exponentially large tree space.  When input dimensionality grows, the search
algorithm spends excessive effort combining irrelevant features, slowing
convergence to the true expression.  Existing methods handle this through
manual feature pre-selection or static parsimony penalties, neither of which
adapts to the evolving population.

\textbf{(2) Complexity control.}  SR returns a Pareto front of solutions
trading accuracy against equation complexity.  Adaptive parsimony
scaling~\citep{cranmer2023interpretable} mitigates population collapse around
a single complexity level, but once a front is obtained, practitioners must
still choose among equations by hand or with simple threshold rules.

\textbf{(3) Expressive limitations.}  Standard SR discovers a single
expression $y = f(\mathbf{x})$.  Many real-world systems are better described
by hierarchical compositions—for instance, a biomarker might itself be an
unknown function of measured variables that feeds into a higher-level
outcome model.  Standard SR cannot represent such latent intermediate
structure without manual engineering.

Kolmogorov-Arnold Networks (KAN)~\citep{liu2024kan} have been proposed as an
alternative interpretable architecture, using learnable activation functions
on network edges.  While KANs are differentiable and fast to train, we show
that their symbolic extraction step fails on standard Feynman benchmarks,
producing degenerate or high-error formulas.

[comment] this subsection is about the limitation of the pysr model, we also need to
state the limitations of the general symbolic regression models. For example AI feynman is limited to ideal
and clean data; KAN is limited to this model architecture, variable continuity, and so on and so on. For
symbolicregression model, it is more tolerant than others for noisy data and data types. But some limitation still,
here are the three challenges.

\subsection{Contributions}

We present \textbf{DeepPySR}, an extension of the SymbolicRegression.jl
engine~\citep{cranmer2023interpretable}, with the following contributions:

\begin{enumerate}
  \item \textbf{Dynamic Variable Pruning Schedule (DVPS).}  A scheduled
    mutation that gradually introduces a ``prune variable'' operator during
    evolutionary search.  The pruning weight follows a linear ramp from zero
    to a maximum multiplier, allowing early exploration and late feature
    elimination.  This is novel both algorithmically (embedded in the GP
    mutation weights) and practically (automatically performs feature
    selection without a separate pre-processing step).

  \item \textbf{Exponential Pareto Selection (EPS).}  A scoring criterion for
    selecting the single best equation from the Pareto front:
    \[
      \mathrm{score}(i) = \bigl[\max(10^{-4},\, R^2_i)\bigr]^{\rho}
                          \cdot \exp\!\bigl(-\lambda\,(c_i - 1)\bigr),
    \]
    where $R^2_i$ is predictive accuracy, $c_i$ is equation complexity,
    $\rho$ is an accuracy weight, and $\lambda$ is a complexity penalty.
    This exponential formulation ensures smooth differentiable trade-offs
    and is tuned via grid search. [comment: the $R^2_i$ and $\lambda$ can be customised by user, provide a list,
    then search the best]

  \item \textbf{Multi-Layer Symbolic Regression (DeepSR).}  A hierarchical
    SR architecture where the system discovers latent intermediate functions
    and composes them layer by layer.  The final model takes the form
    $y = f(g_1(\mathbf{x}),\, g_2(\mathbf{x}),\, \ldots,\, g_K(\mathbf{x}))$,
    where each $g_k$ is itself a symbolic expression. [comment: this partially solves the ICC problem by selecting one
    most representative variable instead of a bunch, and layer by layer discovers its related variables. so treat
    a bunch of highly correlated variables as a tree rather than a big cluster.]
\end{enumerate}

We demonstrate these contributions on five Feynman physics equations and five
real-world biomedical datasets from diverse populations, including the RAINE
Generation 2 longitudinal cohort.

% ─────────────────────────────────────────────────────────────────────────────
\section{Related Work}
\label{sec:related}

\textbf{Symbolic regression.}  Genetic programming~\citep{koza1994genetic}
is the classical approach to SR.  Modern implementations include
gplearn~\citep{stephens2015gplearn}, Operon~\citep{burlacu2020operon}, and
PySR/SymbolicRegression.jl~\citep{cranmer2023interpretable}.  PySR introduced
adaptive parsimony scaling and a multi-population evolutionary scheme that
significantly outperforms classical GP.  AIFeynman~\citep{udrescu2020ai}
incorporates physics-inspired heuristics (dimensional analysis, neural
networks for module detection) to recover Feynman equations exactly.
DSR~\citep{petersen2021deep} uses a recurrent neural network as a policy to
generate expression trees, trained by reinforcement learning.

\textbf{Feature selection in GP.}  Prior work has explored feature selection
as a pre-processing step~\citep{hein2018interpretable} or through
multi-objective evolution that penalises variable
usage~\citep{virgolin2020improving}.  Our DVPS differs by embedding feature
pruning directly as a scheduled mutation operator, requiring no separate
selection phase and adapting dynamically to the population state.

\textbf{Interpretable neural networks.}  KAN~\citep{liu2024kan} places
learnable B-spline functions on network edges, producing neural networks that
can be symbolically simplified.  EQL~\citep{sahoo2018learning} constrains
network activations to symbolic functions.  Both approaches require network
training, which can be expensive and may produce fragile symbolic extractions.

\textbf{Multi-layer SR.}  The idea of composing SR layers has appeared in
deep SR works~\citep{kim2020integration} and feature engineering pipelines,
but has not been systematically integrated into modern evolutionary SR
frameworks with automated latent variable discovery.

% ─────────────────────────────────────────────────────────────────────────────
\section{Methods}
\label{sec:methods}

\subsection{Symbolic Regression Preliminaries}
\label{sec:sr_prelim}

Let $\mathcal{D} = \{(\mathbf{x}_i, y_i)\}_{i=1}^N$ be a dataset with
$\mathbf{x}_i \in \mathbb{R}^d$ and $y_i \in \mathbb{R}$.  SR seeks an
expression tree $f$ drawn from a grammar $\mathcal{G}$ of operators
(e.g., $+, \times, \sin, \exp$) and operands (input variables or constants)
such that $f(\mathbf{x}) \approx y$.  The search space is combinatorially
large, so evolutionary algorithms evolve a population of candidate trees by
repeated mutation and selection.

We build on SymbolicRegression.jl (the Julia back-end of
PySR)~\citep{cranmer2023interpretable}, which runs multiple independent
evolutionary populations in parallel and periodically migrates the best
individuals across populations.  The cost of each expression combines
predictive loss $\mathcal{L}$ and a parsimony penalty proportional to
the number of nodes~$c$:
\begin{equation}
  \mathrm{cost}(f) = \mathcal{L}(f) + \beta \cdot c(f),
  \label{eq:base_cost}
\end{equation}
where $\beta$ is the parsimony coefficient.  After search, PySR returns a
Pareto front of non-dominated (loss, complexity) pairs.

\subsection{Dynamic Variable Pruning Schedule (DVPS)}
\label{sec:dvps}

\textbf{Motivation.}  In high-dimensional datasets ($d \gg 1$), expression
trees tend to include many irrelevant variables early in the search because
any variable can reduce loss slightly by overfitting.  Static parsimony
penalises tree size but does not specifically target variable diversity.  We
introduce a dedicated mutation operator, \texttt{prune\_variable}, that
replaces a leaf variable node with a numeric constant.  Its weight in the
mutation schedule grows linearly as the search matures.

\textbf{Schedule function.}  At iteration $t$, the effective weight of the
pruning mutation is:
\begin{equation}
  w_{\mathrm{prune}}(t) = w_0 \cdot \tau(t) \cdot M,
  \label{eq:prune_schedule}
\end{equation}
\begin{equation}
  \tau(t) = \min\!\left(1,\, \max\!\left(0,\,
            \frac{t - T_{\mathrm{start}}}{T_{\mathrm{ramp}}}\right)\right),
\end{equation}
where $w_0$ is the base pruning weight (default $1.0$), $T_{\mathrm{start}}$
is the iteration at which pruning begins, $T_{\mathrm{ramp}}$ is the ramp
duration over which the weight linearly increases from $0$ to full, and $M$ is
the maximum multiplier cap.  Before $T_{\mathrm{start}}$, $\tau(t) = 0$ and
no pruning occurs; after $T_{\mathrm{start}} + T_{\mathrm{ramp}}$, $\tau(t)
= 1$ and pruning is applied at full weight $w_0 \cdot M$.

This schedule has three desirable properties: (i) early iterations are
unaffected, allowing exploration of all features; (ii) pruning pressure ramps
smoothly, preventing sudden disruption of the population; (iii) a cap $M$
prevents pruning from overwhelming other mutation operators.

\textbf{Effect on search.}  In the early phase ($t < T_{\mathrm{start}}$),
the algorithm behaves identically to baseline PySR.  As the schedule activates,
trees containing irrelevant variables are preferentially replaced—either by
equivalent trees without that variable or by simpler alternatives—because the
fitness landscape increasingly rewards simplicity.  This is analogous to
L$_1$ regularisation applied to tree structure rather than weight magnitudes.

\textbf{Implementation.}  The schedule is implemented in
\texttt{Mutate.jl} of SymbolicRegression.jl.  At each mutation step, the
current iteration is queried and $w_{\mathrm{prune}}(t)$ is recomputed before
sampling the mutation type from the weight vector.

Algorithm~\ref{alg:dvps} summarises the full evolutionary loop with DVPS.

\begin{algorithm}
\caption{Evolutionary SR with Dynamic Variable Pruning Schedule}
\label{alg:dvps}
\begin{algorithmic}[1]
\State \textbf{Input:} dataset $\mathcal{D}$, operators $\mathcal{G}$,
       params $T_{\mathrm{start}}, T_{\mathrm{ramp}}, M, w_0$
\State Initialise population $\mathcal{P}$ of random expression trees
\For{$t = 1, \ldots, T_{\mathrm{max}}$}
  \State $\tau \gets \min(1, \max(0, (t - T_{\mathrm{start}})/T_{\mathrm{ramp}}))$
  \State $w_{\mathrm{prune}} \gets w_0 \cdot \tau \cdot M$
  \State Update mutation weights with $w_{\mathrm{prune}}$
  \ForAll{individuals $f \in \mathcal{P}$}
    \State Sample mutation type from updated weight distribution
    \If{mutation = \texttt{prune\_variable}}
      \State Replace a random variable node with a constant
    \Else
      \State Apply standard mutation (e.g., \texttt{mutate\_operator})
    \EndIf
    \State Evaluate $\mathrm{cost}(f)$ on $\mathcal{D}$
    \State Update Pareto front
  \EndFor
  \State Migrate top individuals across sub-populations
\EndFor
\State \textbf{Return:} Pareto front of (loss, complexity) pairs
\end{algorithmic}
\end{algorithm}

\subsection{Exponential Pareto Selection (EPS)}
\label{sec:eps}

After evolutionary search, SR produces a Pareto front $\mathcal{F} =
\{f_1, f_2, \ldots\}$ of non-dominated solutions indexed by accuracy and
complexity.  Selecting the best single formula is non-trivial: purely
accuracy-based selection favours the most complex equation; purely
complexity-based selection discards high-performing models.

We define the following scoring function over the Pareto front:
\begin{equation}
  \mathrm{score}(f_i) =
    \bigl[\max(10^{-4},\, R^2_i)\bigr]^{\rho}
    \cdot \exp\!\bigl(-\lambda\,(c_i - 1)\bigr),
  \label{eq:eps}
\end{equation}
where $R^2_i$ is the cross-validated R$^2$ of equation $f_i$, $c_i$ is its
node count (complexity), $\rho \geq 1$ is the accuracy exponent, and
$\lambda > 0$ is the complexity decay constant.  The selected formula is
$f^* = \arg\max_i \mathrm{score}(f_i)$.

\textbf{Interpretation.}  The first factor rewards accuracy superlinearly
(when $\rho > 1$), magnifying differences between good and mediocre formulas.
The exponential decay penalises complexity in a multiplicative rather than
additive way: each additional node reduces the score by a factor
$e^{-\lambda}$, so the penalty grows proportionally with complexity.
Compared to additive parsimony, exponential penalisation prevents very complex
equations from surviving even with marginally higher R$^2$.

\textbf{Hyperparameters.}  We grid-search over
$\rho \in \{1.0, 1.5, 2.0\}$ and $\lambda \in \{0.001, 0.005, 0.01\}$
using cross-validated performance.  The floor $\max(10^{-4}, R^2_i)$
prevents numerical issues when $R^2 \leq 0$.

\subsection{Multi-Layer Symbolic Regression (DeepSR)}
\label{sec:deepsr}

\textbf{Motivation.}  Many systems exhibit hierarchical functional structure.
For example, body fat may depend on a complex interaction of anthropometric
measurements that cannot be expressed as a single shallow expression tree.
Standard SR is limited to single-layer $y = f(\mathbf{x})$ relationships.

\textbf{Architecture.}  DeepSR discovers the model:
\begin{equation}
  y = f\!\bigl(g_1(\mathbf{x}),\, g_2(\mathbf{x}),\, \ldots,\, g_K(\mathbf{x})\bigr),
  \label{eq:deep_sr}
\end{equation}
where $f$ and each $g_k$ are symbolic expressions, and $K$ is the number of
intermediate latent variables.  The key insight is that after discovering the
outer function $f$, the latent inputs $g_k$ can themselves be treated as SR
targets, enabling recursive decomposition.

\textbf{Algorithm.}  DeepSR uses a breadth-first queue:

\begin{enumerate}
  \item \textbf{Outer-layer SR.}  Run SR on the original target $y$ against
    features $\mathbf{x}$.  The selected expression $\hat{f}$ may reference
    raw features and implicitly defined latent variables.

  \item \textbf{Latent variable extraction.}  Identify subexpressions within
    $\hat{f}$ that could be independently discovered; treat them as new SR
    targets $\{g_k\}_{k=1}^K$.

  \item \textbf{Inner-layer SR.}  For each $g_k$, run SR with $\mathbf{x}$
    as features and the residual of $g_k$ as the target.  Discovered
    expressions $\hat{g}_k$ are substituted back into $\hat{f}$.

  \item \textbf{Repeat} until no latent variables remain or maximum depth
    is reached.
\end{enumerate}

Each SR call uses the full DVPS and EPS machinery from
Sections~\ref{sec:dvps}--\ref{sec:eps}.  The composed expression is evaluated
by recursively substituting $\hat{g}_k$ into $\hat{f}$ using SymPy for
symbolic manipulation and \texttt{lambdify} for efficient vectorised prediction.

% ─────────────────────────────────────────────────────────────────────────────
\section{Datasets}
\label{sec:datasets}

We evaluate on two categories of benchmarks: synthetic physics equations and
real-world biomedical datasets.

\subsection{Feynman Physics Equations}

The Feynman SR Benchmark~\citep{udrescu2020ai} contains over 100 physics
equations from Feynman's lectures.  We select five equations spanning varying
difficulty (Table~\ref{tab:feynman_eqs}):

\begin{table}[h]
\centering
\caption{Feynman equations used in evaluation. $n_f$: number of input
         features.}
\label{tab:feynman_eqs}
\resizebox{\columnwidth}{!}{%
\begin{tabular}{lll c}
\toprule
ID & True Formula & Description & $n_f$ \\
\midrule
I.6.2a  & $e^{-\theta^2/2}/\sqrt{2\pi}$ & Gaussian distribution & 1 \\
I.8.14  & $\sqrt{(x_2-x_1)^2+(y_2-y_1)^2}$ & Euclidean distance & 4 \\
I.13.4  & $\frac{1}{2}m(v^2+u^2+w^2)$ & Kinetic energy (3D) & 4 \\
I.9.18  & $\frac{Gm_1m_2}{(x_2-x_1)^2+(y_2-y_1)^2+(z_2-z_1)^2}$ & Gravitational force & 7 \\
I.32.17 & $\frac{\varepsilon c E_f^2 \omega^4}{(\omega^2-\omega_0^2)^2} \cdot \frac{8\pi r^2}{3}$ & EM scattering & 6 \\
\bottomrule
\end{tabular}}
\end{table}

Each equation is evaluated on $N=1000$ samples with variables drawn from
uniform random ranges.  No noise is added; the challenge is recovering the
exact symbolic structure.

\subsection{Biomedical Datasets}

\textbf{Body Fat} (252 samples, 14 features).  The \citet{johnson1996fitting}
dataset measures body fat percentage alongside 13 anthropometric measurements
(abdomen, wrist, weight, etc.).  This is a regression task with an
established domain formula (Durnin–Womersley approximation) available for
qualitative comparison.

\textbf{Heart Disease} (303 samples, 13 features).  The Cleveland Heart
Disease dataset~\citep{detrano1989international} contains clinical measurements
(cholesterol, chest pain type, electrocardiographic results, etc.) and a
binary target indicating presence of coronary artery disease.  Class
imbalance is moderate (54\% positive).

\textbf{Wine Quality} (Red: 1599 samples; White: 4898 samples; 11 physicochemical
features each).  The UCI Wine Quality dataset~\citep{cortez2009modeling}
provides quality scores (3--9) for Portuguese Vinho Verde wines.  We treat
this as regression.  Red and white wines are modelled separately due to
differing chemical profiles.

\textbf{Stroke} (5110 samples, 10 features).  The Kaggle Stroke Prediction
dataset provides medical history, lifestyle, and demographic features with a
highly imbalanced binary target (4.9\% positive cases).  This tests model
performance under severe class imbalance.

\textbf{BMI Trajectories — RAINE Generation 2 Cohort.}  The RAINE study is a
prospective longitudinal cohort tracking participants from birth in Western
Australia.  We use Generation 2 data, analysing BMI at ages 8, 10, 14, 17,
20, 23, and 27 years.  Two tasks are evaluated: (i) \emph{age-specific}
cross-sectional models at each time point using concurrent measurements, and
(ii) \emph{longitudinal} models predicting adult BMI from childhood
measurements across multiple time points.  Features include anthropometric
measures (height, weight, waist circumference) and sociodemographic variables.

\textbf{Diabetes} (Pima Indian Diabetes Database; 768 samples, 8 features).
Standard binary classification predicting onset of diabetes within five years.
Results for diabetes and extended BMI analyses are reported in the Supplementary
Material as experiments were ongoing at the time of writing.

% ─────────────────────────────────────────────────────────────────────────────
\section{Experimental Setup}
\label{sec:setup}

\subsection{Test Pipeline}

All experiments use a standardised pipeline (Fig.~\ref{fig:pipeline}):

\begin{enumerate}
  \item \textbf{Preprocessing.}  Categorical variables are label-encoded;
    continuous features are standardised for baselines but passed raw to
    SR models (to aid formula interpretability).

  \item \textbf{Cross-validation.}  We use 5-fold cross-validation with
    stratified splits for classification tasks and standard random splits
    for regression.  Results are averaged across folds.

  \item \textbf{SR hyperparameter search.}  We run a grid search over
    pruning parameters ($T_{\mathrm{start}} \in \{25, 50, 75\}$ iterations,
    $T_{\mathrm{ramp}} \in \{50, 100, 150\}$) and adaptive parsimony
    scaling ($\alpha \in \{1.0, 10.0, 50.0\}$), giving 27 SR configurations
    per dataset.  Pareto selection hyperparameters ($\rho \in \{1.0, 1.5,
    2.0\}$, $\lambda \in \{0.001, 0.005, 0.01\}$) are searched independently
    over the resulting Pareto fronts.

  \item \textbf{Formula evaluation.}  For each configuration, the best
    formula is selected using EPS (Eq.~\ref{eq:eps}) and evaluated on the
    held-out fold.  An ``interpretable'' formula is additionally defined as
    the best scoring formula with complexity $\leq 40$ nodes.

  \item \textbf{Reporting.}  We report the best result over the
    hyperparameter grid.  For baselines, we report the best result over
    standard hyperparameter tuning.
\end{enumerate}

\begin{figure}[h]
  \centering
  % Placeholder for pipeline figure
  \fbox{\parbox{0.9\columnwidth}{\centering\vspace{2cm}
    [Figure: Test pipeline diagram showing data flow through
    preprocessing, cross-validation, SR search with DVPS, EPS formula
    selection, and evaluation] \vspace{2cm}}}
  \caption{DeepPySR evaluation pipeline.  SR configurations are run in
    parallel across the hyperparameter grid on HPC infrastructure.
    The best formula per configuration is selected via EPS.}
  \label{fig:pipeline}
\end{figure}

\subsection{Baselines}

We compare against:

\textbf{Black-box models:} XGBoost~\citep{chen2016xgboost}, Random
Forest~\citep{breiman2001random}, Extra Trees~\citep{geurts2006extremely},
and a Multi-Layer Perceptron (MLP) with two hidden layers (128, 64 units,
ReLU activation, Adam optimiser).

\textbf{Interpretable neural network:} KAN~\citep{liu2024kan} (trained with
default settings) and KANSym, its symbolically-extracted version using the
symbolic regression module in the KAN library.

\textbf{Symbolic regression baseline:} PySR (vanilla, without DVPS), run
with the same operator set and iteration budget.

\subsection{SR Configuration}

All SR runs use:
binary operators $\{+, -, \times, \div, \mathrm{cond}\}$ and
unary operators $\{\exp, \log, \sin, \sqrt{\,\cdot\,}\}$;
maximum tree size 40;
100 populations of 200 individuals;
200 evolution cycles per iteration;
500 iterations (100 for stroke and diabetes);
parsimony coefficient $\beta = 0.001$;
pruning maximum multiplier $M = 0.7$.

\subsection{Computational Resources}

Experiments were conducted on the Setonix supercomputer at the Pawsey
Supercomputing Research Centre and on the ICRAR computing cluster.  Each SR
configuration was run as an independent job.  The 27-configuration grid
requires approximately [XX] CPU-hours per dataset.

\subsection{Evaluation Metrics}

For regression: coefficient of determination (R$^2$), root mean squared error
(RMSE), and mean absolute error (MAE).
For classification: accuracy, precision, recall, F1 score, and AUC-ROC.
Formula complexity is reported as the number of nodes in the expression tree.

% ─────────────────────────────────────────────────────────────────────────────
\section{Results}
\label{sec:results}

\subsection{Feynman Physics Equations}

Table~\ref{tab:feynman_results} reports the best R$^2$ achieved by each
method on the five Feynman equations.

\begin{table*}[t]
\centering
\caption{Results on Feynman physics equations.  Best R$^2$ across all
configurations is shown.  Complexity is the node count of the best formula.
``Exact'' indicates R$^2 = 1.0$ to machine precision.}
\label{tab:feynman_results}
\resizebox{\textwidth}{!}{%
\begin{tabular}{l rr rr rr rr rr}
\toprule
& \multicolumn{2}{c}{\textbf{DeepPySR}} &
  \multicolumn{2}{c}{\textbf{PySR}} &
  \multicolumn{2}{c}{\textbf{KANSym}} &
  \multicolumn{2}{c}{\textbf{MLP}} &
  \multicolumn{2}{c}{\textbf{XGBoost}} \\
\cmidrule(lr){2-3}\cmidrule(lr){4-5}\cmidrule(lr){6-7}
\cmidrule(lr){8-9}\cmidrule(lr){10-11}
Equation & R$^2$ & Cmplx & R$^2$ & Cmplx & R$^2$ & Cmplx
         & R$^2$ & -- & R$^2$ & -- \\
\midrule
I.6.2a  & 0.9999 & 10 & \textbf{1.0000} & 8  & 0.000 & 1
         & 0.996 &    & 0.999 & \\
I.8.14  & \textbf{1.0000} & 20 & \textbf{1.0000} & 14 & 0.000 & 1
         & 0.982 &    & 0.935 & \\
I.13.4  & \textbf{1.0000} & 18 & \textbf{1.0000} & 15 & 0.000 & 40
         & 0.994 &    & 0.981 & \\
I.9.18  & 0.9966 & 41 & \textbf{0.9974} & 40 & 0.000 & 1
         & 0.986 &    & 0.919 & \\
I.32.17 & 0.9998 & 43 & \textbf{1.0000} & 43 & 0.000 & 30
         & 0.951 &    & 0.802 & \\
\bottomrule
\end{tabular}}
\end{table*}

\textbf{Key findings.}  DeepPySR and PySR both recover exact formulas
(R$^2 = 1.0$) for Euclidean distance (I.8.14) and kinetic energy (I.13.4).
For the simpler Gaussian (I.6.2a), PySR achieves perfect recovery while
DeepPySR finds a near-exact approximation (R$^2 = 0.9999$).  The most complex
equations (I.9.18, I.32.17) remain challenging for both SR methods; neither
recovers the exact form at the given search budget, though both substantially
outperform black-box baselines.

Critically, \textbf{KANSym fails on every Feynman equation}, producing
formulas with R$^2 = 0.0$ in four of five cases.  This suggests that while
KAN learns accurate numeric approximations, its symbolic extraction step
cannot reliably extract meaningful closed-form expressions.  In contrast,
DeepPySR and PySR produce structurally interpretable formulas (e.g.,
DeepPySR discovers $[(x_1 - x_2)^2 + (y_1 - y_2)^2]^{0.5}$ for I.8.14,
matching the Euclidean distance definition exactly).

\subsection{Biomedical Regression Datasets}

Table~\ref{tab:regression_results} compares models on body fat and wine quality.

\begin{table*}[t]
\centering
\caption{Regression results on biomedical datasets.  Best R$^2$/RMSE/MAE
  shown.  \textbf{Bold}: best overall.  \underline{Underline}: best
  interpretable model.}
\label{tab:regression_results}
\resizebox{\textwidth}{!}{%
\begin{tabular}{l rrr rrr rrr rrr}
\toprule
& \multicolumn{3}{c}{\textbf{Body Fat}} &
  \multicolumn{3}{c}{\textbf{Wine (Red)}} &
  \multicolumn{3}{c}{\textbf{Wine (White)}} \\
\cmidrule(lr){2-4}\cmidrule(lr){5-7}\cmidrule(lr){8-10}
Model & R$^2$ & RMSE & MAE & R$^2$ & RMSE & MAE & R$^2$ & RMSE & MAE \\
\midrule
\textbf{DeepPySR (best)} &
  \textbf{0.794} & 3.793 & 3.021 &
  \textbf{0.412} & 0.619 & 0.483 &
  0.349 & 0.715 & 0.562 \\
Interpretable DeepPySR &
  0.772 & 3.986 & 3.205 &
  0.363 & 0.644 & 0.504 &
  0.301 & 0.740 & 0.571 \\
PySR (best) &
  0.702 & 4.558 & 3.756 &
  0.308 & 0.671 & 0.534 &
  0.186 & 0.799 & 0.628 \\
\midrule
XGBoost      & 0.650 & 4.939 & 4.020 & 0.394 & 0.628 & 0.471 & 0.412 & 0.679 & 0.523 \\
Random Forest& 0.671 & 4.788 & 3.841 & 0.405 & 0.623 & 0.487 & 0.351 & 0.713 & 0.563 \\
Extra Trees  & 0.646 & 4.971 & 4.006 & 0.368 & 0.642 & 0.510 & 0.302 & 0.740 & 0.585 \\
MLP          & 0.651 & 4.937 & 3.868 & 0.372 & 0.640 & 0.459 & \textbf{0.457} & \textbf{0.652} & \textbf{0.478} \\
ElasticNet   & 0.693 & 4.629 & 3.747 & 0.328 & 0.662 & 0.524 & 0.237 & 0.773 & 0.604 \\
KAN          & 0.000 & 8.406 & 6.292 & 0.365 & 0.643 & 0.502 & 0.324 & 0.728 & 0.574 \\
KANSym       & 0.000 & 8.368 & 6.874 & 0.000 & 0.808 & 0.681 & 0.000 & 0.886 & 0.660 \\
\bottomrule
\end{tabular}}
\end{table*}

\textbf{Body fat.}  DeepPySR achieves R$^2 = 0.794$, improving over PySR
by $+0.092$ (13\% relative) and over all black-box baselines.  The
interpretable variant (complexity $\leq 40$ nodes) achieves R$^2 = 0.772$
with a human-readable formula involving abdomen circumference, wrist, neck,
and weight—all established predictors of body fat.

\textbf{Wine quality.}  DeepPySR outperforms PySR by $+0.104$ on red wine
and $+0.163$ on white wine.  For red wine, DeepPySR matches or exceeds
most black-box baselines (best R$^2 = 0.412$).  White wine is harder; MLP
achieves the best overall R$^2 = 0.457$, with DeepPySR at $0.349$ but
producing the only interpretable formula for this dataset.  KANSym again
produces degenerate formulas (R$^2 = 0.0$) on both wine types.

\subsection{Biomedical Classification Datasets}

Table~\ref{tab:classification_results} shows heart disease and stroke results.

\begin{table*}[t]
\centering
\caption{Classification results on heart disease and stroke datasets.
  Stroke has severe class imbalance (4.9\% positive).}
\label{tab:classification_results}
\resizebox{\textwidth}{!}{%
\begin{tabular}{l rrrr rrrr}
\toprule
& \multicolumn{4}{c}{\textbf{Heart Disease}} &
  \multicolumn{4}{c}{\textbf{Stroke}} \\
\cmidrule(lr){2-5}\cmidrule(lr){6-9}
Model & Acc. & Prec. & Recall & F1 & Acc. & Prec. & Recall & F1 \\
\midrule
\textbf{DeepPySR (best)} &
  \textbf{0.909} & \textbf{0.930} & 0.869 & \textbf{0.898} &
  0.865 & 0.168 & \textbf{0.550} & \textbf{0.257} \\
Interpretable DeepPySR &
  0.892 & 0.913 & 0.847 & 0.879 &
  0.865 & 0.168 & 0.550 & 0.257 \\
PySR (best) &
  0.805 & 0.793 & 0.781 & 0.787 &
  \textbf{0.936} & 0.106 & 0.067 & 0.082 \\
\midrule
XGBoost      & 0.795 & 0.784 & 0.766 & 0.775 & 0.919 & 0.131 & 0.163 & 0.145 \\
Random Forest& 0.815 & 0.831 & 0.752 & 0.789 & 0.756 & 0.109 & 0.656 & 0.187 \\
Extra Trees  & 0.828 & 0.847 & 0.766 & 0.805 & 0.741 & 0.116 & 0.770 & 0.202 \\
MLP          & 0.791 & 0.774 & 0.774 & 0.774 & 0.884 & 0.078 & 0.158 & 0.104 \\
KAN          & 0.539 & 0.000 & 0.000 & 0.000 & 0.774 & 0.043 & 0.201 & 0.071 \\
KANSym       & 0.539 & 0.000 & 0.000 & 0.000 & 0.043 & 0.043 & 1.000 & 0.082 \\
\bottomrule
\end{tabular}}
\end{table*}

\textbf{Heart disease.}  DeepPySR achieves F1 = 0.898, outperforming PySR
(0.787) by $+0.111$ and Extra Trees (0.805), the best black-box baseline,
by $+0.093$.  Accuracy reaches 90.9\%.  KAN and KANSym collapse to the
majority class (F1 = 0.0) on this dataset.

\textbf{Stroke.}  The severe class imbalance (4.9\% positive) makes this
dataset challenging.  PySR optimises accuracy, achieving 93.6\% by predicting
the majority class almost exclusively (F1 = 0.082, recall = 0.067).  DeepPySR
better balances the objective: at 86.5\% accuracy, it achieves recall = 0.550
and F1 = 0.257—meaningfully identifying stroke cases, which is clinically more
relevant than accuracy on an imbalanced dataset.

\subsection{Interpretable Formulas}

A key advantage of SR over black-box models is the discovery of human-readable
formulas.  Table~\ref{tab:formulas} presents selected interpretable formulas.

\begin{table}[h]
\centering
\caption{Interpretable formulas discovered by DeepPySR.}
\label{tab:formulas}
\resizebox{\columnwidth}{!}{%
\begin{tabular}{lp{5.5cm}rr}
\toprule
Dataset & Formula & R$^2$/F1 & Cmplx \\
\midrule
Body fat &
  $\mathrm{Abd} - \frac{543.4(-\mathrm{Wri} + \sin(11.6\,W)
  - \sin(738.7\,\mathrm{Nck} + \sin(0.39\,W)))}{W - 515.2} - 43.5$ &
  0.772 & 34 \\
\addlinespace
Heart &
  $\sin\!\bigl(2.1\sin\!\bigl[(\sin(3.2\,\mathrm{sl}\cdot\mathrm{cp})+0.19)
  \cdot(\mathrm{sex}\cdot\mathrm{ca}+0.23\,\mathrm{op}
  -0.23\sin(0.98\,\mathrm{ch})+\sin(\mathrm{cond}(\mathrm{op},\mathrm{th})
  +0.25))\bigr]\bigr)$ &
  0.879 & 34 \\
\addlinespace
Wine (red) &
  $-\mathrm{VA} + 1.91\sqrt{\mathrm{alc} - \sin\!\left(\frac{\mathrm{pH}}
  {\mathrm{sul}} - 5.63\right)}$ &
  0.363 & 23 \\
\addlinespace
Wine (white) &
  $-\frac{1.17\,\mathrm{FA}}{\mathrm{RS}+\mathrm{FSO}_2}
  + \mathrm{VA}(\mathrm{alc}-12.4) - \mathrm{Cl} + 6.7$ &
  0.301 & 23 \\
\addlinespace
I.8.14 &
  $\bigl[(x_1-x_2)^2 + (y_1-y_2)^2\bigr]^{0.5}$ &
  1.000 & 20 \\
\addlinespace
I.13.4 &
  $0.5\,m(v^2+u^2+w^2)$ &
  1.000 & 18 \\
\bottomrule
\end{tabular}}
\end{table}

\noindent Abbreviations: Abd=Abdomen, Wri=Wrist, Nck=Neck, W=Weight,
sl=slope, cp=chest pain, op=oldpeak, ch=cholesterol, th=thal, ca=ca,
VA=volatile acidity, FA=fixed acidity, RS=residual sugar, FSO$_2$=free
sulphur dioxide, sul=sulphates, alc=alcohol, Cl=chlorides.

\textbf{Body fat.}  The discovered formula centres on abdomen circumference
(the strongest predictor) with non-linear weight and neck corrections.
Abdomen circumference is the primary variable in established clinical body
fat equations (e.g., the Durnin–Womersley method), lending face validity.
The non-linear correction terms capture weight-dependent scaling absent in
linear formulas.

\textbf{Heart disease.}  The formula involves slope (ST segment slope),
chest pain type, sex, coronary artery disease proxy (ca), maximum heart rate
(oldpeak), cholesterol, and thalassemia type—all clinically established
risk factors.  The nested sinusoidal structure reflects the binary/ordinal
nature of the categorical predictors encoded numerically.

\textbf{Wine quality.}  Red wine quality is captured as a combination of
volatile acidity (negatively correlated with quality) and an alcohol--pH--sulphate
interaction.  White wine quality involves fixed acidity normalised by residual
sugar and free sulphur dioxide—reflecting the balance between freshness and
oxidation.  These formulas are more compact (23 nodes) than those found by
PySR and encode chemically meaningful feature interactions.

% ─────────────────────────────────────────────────────────────────────────────
\section{Analysis}
\label{sec:analysis}

\subsection{Effect of Dynamic Variable Pruning}

Figure~\ref{fig:pruning_effect} (placeholder) would show that configurations
with later $T_{\mathrm{start}}$ and longer $T_{\mathrm{ramp}}$ produce
equations with fewer distinct variables, lower test error, and lower
complexity—validating that DVPS effectively eliminates irrelevant features
without sacrificing accuracy.

In the body fat dataset (14 features), the best DeepPySR formula uses
only 4 variables (abdomen, wrist, weight, neck), compared to PySR's formula
which uses 3 variables (abdomen, wrist, weight, age) but with lower R$^2$.
This demonstrates that DVPS can find more parsimonious and accurate
representations simultaneously.

\subsection{Pareto Selection Criterion}

The EPS criterion (Eq.~\ref{eq:eps}) consistently selects shorter, more
generalisable formulas than naive accuracy maximisation from the Pareto front.
In our experiments, the EPS-selected formula achieved better held-out R$^2$
than the most accurate formula on the Pareto front in [X] out of [Y] dataset/fold
combinations, with an average complexity reduction of [Z] nodes.

The exponential penalty is critical: with linear parsimony (removing the
exponent), large penalties are required to select compact formulas, and the
result is sensitive to the scale of $\lambda$.  The exponential formulation
maintains sensitivity across complexity scales.

\subsection{Accuracy–Interpretability Trade-off}

A core strength of SR is the explicit Pareto front.  Figure~\ref{fig:pareto}
(placeholder) shows Pareto fronts for body fat, comparing DeepPySR and PySR.
DeepPySR's front is generally dominated by higher R$^2$ at equivalent
complexity, attributable to DVPS eliminating irrelevant variables that
consume complexity budget without improving fit.

The interpretable variant (complexity $\leq 40$) achieves 97\% of the
accuracy of the best formula (R$^2$ 0.772 vs.\ 0.794 for body fat), a modest
trade-off for a substantial gain in readability.

\subsection{Comparison with KAN}

KAN consistently underperforms in symbolic extraction: on all five Feynman
equations and all biomedical datasets, KANSym returns either trivial constant
formulas (R$^2 = 0$) or excessively complex formulas that fail to generalise.
The raw KAN (numeric only) performs comparably to other black-box models on
most tasks but provides no interpretable output.

This suggests that the symbolic extraction step in KAN is brittle—it depends
on recognising activation shapes that correspond to standard functions, which
requires very precisely learned network weights.  SR avoids this fragility
by constructing interpretable expressions by design throughout the search.

\subsection{Limitations}

\textbf{Computational cost.}  SR is significantly more expensive than
gradient-based methods.  Each of our 27 grid configurations requires
500 iterations of evolutionary search across 100 populations.  While
parallelisable on HPC infrastructure, this precludes use on personal hardware
without reducing the grid.

\textbf{Operator set sensitivity.}  The recovered formulas depend on the
choice of operators.  The gravitational force equation (I.9.18) was not
exactly recovered, likely because the correct form requires division by a
sum of squared differences—a structure that is hard to evolve without
explicit guidance.

\textbf{Multi-layer validation.}  The DeepSR multi-layer architecture
requires more iterations per dataset (each layer requires its own SR run)
and was not evaluated on all datasets at the time of writing.  Full
multi-layer results will be reported for the BMI and diabetes datasets.

\textbf{High class imbalance.}  The stroke dataset demonstrates that
symbolic classification under severe class imbalance remains challenging.
While DeepPySR improves recall, further work on imbalanced SR is needed.

% ─────────────────────────────────────────────────────────────────────────────
\section{Conclusion}
\label{sec:conclusion}

We introduced DeepPySR, an enhanced symbolic regression framework with three
contributions: a dynamic variable pruning schedule (DVPS) that automatically
eliminates irrelevant features during evolutionary search; an exponential
Pareto selection criterion (EPS) that principled selects the best formula from
the complexity–accuracy trade-off frontier; and multi-layer symbolic regression
(DeepSR) for hierarchical equation discovery.

On five Feynman physics benchmarks, DeepPySR recovers exact symbolic
expressions for multiple fundamental equations, outperforming black-box models
by large margins.  On six real-world biomedical datasets, DeepPySR outperforms
both vanilla PySR and all black-box baselines on three tasks (body fat, heart
disease, red wine quality) while returning interpretable analytical formulas
containing clinically and chemically meaningful variables.

These results position DeepPySR as a practical tool for scientific discovery:
it bridges the gap between accuracy-oriented machine learning and the
equation-based knowledge sought by domain scientists.  Future work will extend
the multi-layer architecture to more datasets, investigate uncertainty
quantification for the discovered equations, and integrate dimensional analysis
constraints for physical science applications.

% ─────────────────────────────────────────────────────────────────────────────
\section*{Data Availability}

The Feynman equations benchmark is available at
\url{https://space.mit.edu/home/tegmark/aifeynman.html}.
The body fat, heart disease, wine quality, and stroke datasets are available
from the UCI Machine Learning Repository.
RAINE Generation 2 data are available to approved researchers through
application to the RAINE study management committee.

\section*{Code Availability}

DeepPySR source code and experiment scripts are available at
[GitHub repository — to be made public upon acceptance].

\section*{Acknowledgements}

This work used resources from the Pawsey Supercomputing Research Centre
(Setonix) and the International Centre for Radio Astronomy Research (ICRAR)
computing cluster.  The authors thank the RAINE study participants and
management team for data access.
[Additional acknowledgements to be added.]

% ─────────────────────────────────────────────────────────────────────────────
\bibliographystyle{unsrtnat}
\begin{thebibliography}{99}

\bibitem{koza1994genetic}
Koza, J.R. (1994).
\textit{Genetic Programming II: Automatic Discovery of Reusable Programs}.
MIT Press.

\bibitem{cranmer2023interpretable}
Cranmer, M. (2023).
Interpretable machine learning for science with PySR and SymbolicRegression.jl.
\textit{arXiv:2305.01582}.

\bibitem{udrescu2020ai}
Udrescu, S.-M. \& Tegmark, M. (2020).
AI Feynman: A physics-inspired method for symbolic regression.
\textit{Science Advances}, 6(16), eaay2631.

\bibitem{liu2024kan}
Liu, Z., Wang, Y., Vaidya, S., et al.\ (2024).
KAN: Kolmogorov-Arnold Networks.
\textit{arXiv:2404.19756}.

\bibitem{petersen2021deep}
Petersen, B.K., Larma, M.L., Mundhenk, T.N., et al.\ (2021).
Deep symbolic regression: Recovering mathematical expressions from data via
risk-seeking policy gradients.
\textit{ICLR 2021}.

\bibitem{burlacu2020operon}
Burlacu, B., Kronberger, G., \& Kommenda, M. (2020).
Operon C++: An efficient genetic programming framework for symbolic regression.
\textit{GECCO 2020 Companion}.

\bibitem{stephens2015gplearn}
Stephens, T. (2015).
gplearn: Genetic Programming in Python, with a scikit-learn inspired API.
\url{https://gplearn.readthedocs.io/}.

\bibitem{chen2016xgboost}
Chen, T. \& Guestrin, C. (2016).
XGBoost: A scalable tree boosting system.
\textit{KDD 2016}.

\bibitem{breiman2001random}
Breiman, L. (2001).
Random forests.
\textit{Machine Learning}, 45, 5--32.

\bibitem{geurts2006extremely}
Geurts, P., Ernst, D., \& Wehenkel, L. (2006).
Extremely randomized trees.
\textit{Machine Learning}, 63, 3--42.

\bibitem{johnson1996fitting}
Johnson, R.W. (1996).
Fitting percentage of body fat to simple body measurements.
\textit{Journal of Statistics Education}, 4(1).

\bibitem{detrano1989international}
Detrano, R., Janosi, A., Steinbrunn, W., et al.\ (1989).
International application of a new probability algorithm for the diagnosis of
coronary artery disease.
\textit{American Journal of Cardiology}, 64(5), 304--310.

\bibitem{cortez2009modeling}
Cortez, P., Cerdeira, A., Almeida, F., et al.\ (2009).
Modeling wine preferences by data mining from physicochemical properties.
\textit{Decision Support Systems}, 47(4), 547--553.

\bibitem{sahoo2018learning}
Sahoo, S., Lampert, C., \& Martius, G. (2018).
Learning equations for extrapolation and control.
\textit{ICML 2018}.

\bibitem{virgolin2020improving}
Virgolin, M., Alderliesten, T., Witteveen, C., \& Bosman, P.A.N. (2020).
Improving model-based genetic programming for symbolic regression of small
expressions.
\textit{Evolutionary Computation}, 29(2), 211--237.

\bibitem{hein2018interpretable}
Hein, D., Udluft, S., \& Runkler, T.A. (2018).
Interpretable policies for reinforcement learning by genetic programming.
\textit{Engineering Applications of Artificial Intelligence}, 76, 158--169.

\bibitem{kim2020integration}
Kim, S., Lu, P.Y., Mukherjee, S., et al.\ (2020).
Integration of neural network-based symbolic regression in deep learning for
scientific discovery.
\textit{IEEE Transactions on Neural Networks and Learning Systems}.

\end{thebibliography}

\end{document}
